\section{Definição do Problema de Controle}
\label{sec:definicao_problema_controle}
O controle é utilizado para exercer as ações de propulsão e de atitude do corpo, bem como os torques articulares do manipulador, de modo que:
(i) na aproximação final, a orientação relativa satisfaça tolerâncias angulares definidas para captura;
(ii) limites cinemáticos e dinâmicos sejam respeitados;
(iii) a reação na base seja compensada de forma a atender os limites estabelecidos.
A técnica primária adotada neste estudo é o controle proporcional--derivativo (PD) aplicado em múltiplas malhas, com complementos de compensação de modelo, filtragem do termo derivativo e saturação quando aplicável.

\subsection{Estados e Variáveis de Controle}
\label{sec:var_estado_controle}

Para este trabalho, adotou-se um conjunto de variáveis de estado, conforme:
\begin{equation}
    \label{eq:variaveis_estado}
var_{estado} = \{ \vec r,\; \vec v,\; q,\; \vec \omega,\; \theta,\; \dot\theta \},
\end{equation}
onde $\vec r$ e $\vec v$ representam posição e velocidade relativas no referencial LVLH, $q$ é o quaternion unitário da base, $\vec \omega$ a velocidade angular corporal, $\theta$ os ângulos articulares e $\dot\theta$ suas velocidades.
Por sua vez, as variáveis de controle utilizadas são:
\begin{equation}
    \label{eq:variaveis_controle}
F_{controle} = \{ \vec f,\; \vec{\tau}\},
\end{equation}
em que $\vec f$ são forças, $\vec{\tau}$ são os torques.
O modelo dinâmico inclui as equações orbitais, sendo o modelo do problema de dois corpos, as Equações de Hill/CW linearizadas, as equações de atitude do corpo rígido e as equações do manipulador robótico.


\subsection{Restrições dos Atuadores}

Além dos limites cinemáticos e geométricos das juntas e dos elos do manipulador robótico (MR), ele está sujeito às restrições dinâmicas implementadas nos atuadores, que são expressas em termos de torque máximo nas juntas revolutas e força máxima na junta prismática.  
As restrições são expressas conforme:
\begin{equation}
\label{eq:limite_torque}
|\tau_j| \leq \tau_{j,max} \qquad \text{(juntas revolutas)},
\end{equation}

\begin{equation}
\label{eq:limite_forca}
|f_5| \leq f_{d5,max} \qquad \text{(junta prismática)}.
\end{equation}

% \begin{equation}
% \label{eq:limit_forca}
% |f_5(t)| \leq f_{5,\max}, \qquad
% |\tau_i(t)| \leq \tau_{i,\max}, \;\; i=1,\dots,4.
% \end{equation}

Esses valores máximos foram definidos de forma arbitrária, utilizando como base as simulações realizadas.
% , e os valores de $\tau_{j,\text{max}}$ e $f_{d5,max}$ serão apresentados posteriormente com base na análise dinâmica completa e na sintonia do controlador implementado no ambiente de simulações \emph{MATLAB/Simulink}.  


% \subsection{Critérios de Desempenho}
\label{sec:criterios_desempenho}

Os critérios de desempenho permitem quantificar indicares de qualidade do modelo implementado. As métricas utilizadas são:
\begin{itemize}
  \item erro (desejado e atual) de posição e de ângulo;
  \item tempo de estabilização;
  \item esforço do controle, onde é avaliado a existência de torques ou de forças irreais;
  \item comportamento em regime permanente;
  \item saturação dos atuadores, verificando se os torques ou forças exigidos permanecem dentro dos limites atribuídos ao sistema;
\end{itemize}

As métricas completas podem ser visualizadas no Apêndice \eqref{ApendiceA}.




